\documentclass[12pt]{article}

% --- 1. Core Packages ---
\usepackage[utf8]{inputenc}
\usepackage[vietnamese, english]{babel} % Supports both if needed
\usepackage[margin=1in]{geometry}       % Standard Word-style margins
\usepackage{amsmath, amssymb}            % Essential for Data Science math
\usepackage{graphicx}                   % For including images
\usepackage{booktabs}                   % For professional-looking tables
\usepackage{xcolor}                     % for code coloring

% --- 2. Code Block Setup (listings) ---
\usepackage{listings}
\definecolor{codegreen}{rgb}{0,0.6,0}
\lstset{
    basicstyle=\ttfamily\small,
    commentstyle=\color{codegreen},
    keywordstyle=\color{blue},
    numbers=left,
    numberstyle=\tiny\color{gray},
    frame=single,
    breaklines=true
}

% --- 3. Document Metadata ---
\title{Technical Report: Multimodal Analysis and RAG Injection}
\author{Bao --- Student ID: [Your ID]}
\date{\today}

\begin{document}

\maketitle % Generates the title at the top

\section{Introduction}
This is a placeholder for your introduction. Since you are working on Image Editing and RAG, you might describe your environment here. LaTeX is excellent for maintaining structure in long documents.

\section{Mathematical Foundation}
In Data Science, we often need to define loss functions or probability distributions. Here is an example of a squared error loss with regularization:

\begin{equation}
    \mathcal{L}(\theta) = \sum_{i=1}^{n} (y_i - f(x_i; \theta))^2 + \lambda \|\theta\|^2
\end{equation}

\section{Implementation}
You can include your Python code directly in the document. This is much cleaner than taking screenshots of your IDE:

\begin{lstlisting}[language=Python, caption=Simple Data Preprocessing]
import numpy as np

def normalize_data(data):
    # Scale data between 0 and 1
    return (data - np.min(data)) / (np.max(data) - np.min(data))

print("Data normalized successfully!")
\end{lstlisting}

\section{Experimental Results}
Below is how you would typically include a result image from your Image Editing experiments:

\begin{figure}[h]
    \centering
    %\includegraphics[width=0.7\textwidth]{result_plot.png} 
    \small (Image Placeholder: Uncomment the line above once you have a real image file in this folder)
    \caption{Output of the Spectral RAG Injection process.}
\end{figure}

\section{Conclusion}
Everything looks ready for your thesis work. Happy TeX-ing!

\end{document}